\chapter{Despliegue y Puesta en Producción}
\label{cap:despliegue}

Una etapa fundamental en el ciclo de vida del desarrollo de software es la transición del entorno de desarrollo local al entorno de producción. En este capítulo se documentan las decisiones técnicas, la configuración del servidor y los procesos realizados para que el sistema de \textbf{Turnos Académicos} del Observatorio Alanis quede accesible a la comunidad universitaria, así como el proceso de transferencia del sistema a los usuarios finales.

\section{Preparación para Producción}

Antes de trasladar el código al servidor, fue necesario realizar ajustes de optimización y seguridad propios del entorno de producción. Entre las tareas destacadas se encuentran:

\begin{itemize}
    \item \textbf{Gestión de dependencias:} Instalación exclusiva de paquetes requeridos para producción, omitiendo herramientas de desarrollo (ejecutando \texttt{composer install --optimize-autoloader --no-dev}).
    \item \textbf{Compilación de Assets:} Los archivos de \textit{frontend} y componentes de Vue.js fueron compilados y minimizados para optimizar los tiempos de carga en los navegadores de los usuarios.
    \item \textbf{Caché de configuración:} Se generaron cachés de rutas, vistas y configuración de Laravel (\texttt{config:cache}, \texttt{route:cache}) para acelerar los tiempos de respuesta del framework.
    \item \textbf{Ajustes de Seguridad (\texttt{.env}):} Desactivación del modo depuración (\texttt{APP\_DEBUG=false}), configuración de la clave de encriptación de la aplicación (\texttt{APP\_KEY}) y establecimiento de las credenciales de conexión a la base de datos de producción.
\end{itemize}

\section{Entorno de Servidor y Despliegue (Deploy)}

El despliegue de la aplicación se llevó a cabo en la infraestructura provista por la Facultad de Ciencias Exactas de la Universidad Nacional de Salta (\textbf{UNSa}), asignada al subdominio \texttt{observatorioalanis.exa.unsa.edu.ar}.

La configuración del entorno y el procedimiento de despliegue consistió en:

\begin{enumerate}
    \item \textbf{Infraestructura y Hosting Institucional:} El sistema se aloja en el entorno de hosting compartido administrado por la Supervisión de RED de la facultad mediante la plataforma \textbf{Usermin / cPanel}.
    \item \textbf{Transferencia de Archivos (FTP):} Se utilizó un cliente FTP (\textit{FileZilla}) para transferir los archivos de la aplicación a la carpeta asignada (\texttt{public\_html/}), configurando el directorio raíz del servidor web hacia el directorio \texttt{public} de Laravel para proteger la estructura del framework y los archivos de configuración.
    \item \textbf{Servidor de Base de Datos (MariaDB):} La persistencia de datos en producción se gestiona sobre un servidor **MariaDB** proporcionado por la facultad. El esquema relacional y la carga inicial de tablas (dependencias, trámites y estructura de usuarios) se realizaron mediante la importación del archivo estructurado de base de datos desde el módulo MySQL de Usermin.
    \item \textbf{Configuración del Servidor de Correo (SMTP):} Se establecieron los parámetros SMTP necesarios para el envío automático de comprobantes y notificaciones de confirmación de turnos a los correos de los usuarios.
\end{enumerate}

\section{Capacitación y Entrega al Personal}

Para garantizar una transición fluida y el correcto aprovechamiento del sistema por parte del personal del Observatorio Alanis y de la universidad, se realizó un proceso de transferencia de conocimientos y material de apoyo.

\subsection{Material de Soporte y Manuales}

Se estructuró una guía de usuario para los administradores y operadores del sistema, contemplando:
\begin{itemize}
    \item Alta y configuración de dependencias, trámites y franjas horarias de atención.
    \item Gestión del calendario de días no laborables y feriados institucionales.
    \item Procedimiento para la atención presencial y actualización del estado de los turnos.
    \item Generación e impresión de reportes de asistencia y estadísticas de visitas al Observatorio.
\end{itemize}

\subsection{Sesiones de Capacitación}

Se llevaron a cabo jornadas de capacitación dirigidas al personal administrativo encargados de la gestión de agenda del Observatorio Alanis. Durante las sesiones se realizó un recorrido funcional completo por el panel administrativo, capacitando al personal en el control de capacidad de atención ("mesas habilitadas") y resolución de incidencias en la asignación de turnos.